\section{Troubleshooting}
\label{sec:trouble}

The entries below are ordered by the stage at which the message appears. Each quotes the
characteristic part of the message, in the form in which PowerShell prints it. Red text is
PowerShell's way of marking an error; read its first line, which names the cause, and ignore
the lines beginning with a plus sign that follow it.

Several remedies ask for the \emph{opening sequence}. It consists of the following two
commands, given in this order in a PowerShell window; afterwards the prompt begins with
\texttt{(.venv)}.
\begin{ps}[Opening sequence, command 1 of 2]
cd C:\ee105\camera-to-latex
\end{ps}
\begin{ps}[Opening sequence, command 2 of 2]
.\.venv\Scripts\Activate.ps1
\end{ps}

\subsection{PowerShell and the tools}

\trouble{\texttt{brew : The term 'brew' is not recognized as the name of a cmdlet}}
\cause \cmd{brew} is the package manager of macOS. Instructions found on the internet often assume a
Mac.
\remedy Use only the commands of this guide; on Windows the corresponding program is \cmd{winget}.

\trouble{\texttt{winget : The term 'winget' is not recognized}}
\cause The Windows component \emph{App Installer} is missing or out of date.
\remedy Open Microsoft Store, search for \emph{App Installer}, click \ui{Get} or \ui{Update}, and
open a new PowerShell window.

\trouble{\texttt{git} (or \texttt{py}, \texttt{pdflatex}, \texttt{pdftoppm}) \texttt{: The
term \ldots{} is not recognized}}
\cause Either the window was opened before the program was installed, or the installation failed.
\remedy Open a new PowerShell window and try again. If the message persists, repeat the installation
step of that program in Section~\ref{sec:tools} and read its output for an error.

\trouble{\texttt{pdftoppm} is still not recognized after Poppler has been installed again}
\cause The folder that contains \file{pdftoppm.exe} is not on the list of folders that Windows
searches (the \term{path}).
\remedy Run the five commands below in order. The third must print the name of a folder that ends in
\file{Library\bs bin}; if it prints nothing, Poppler is not installed, and the remedy is to
repeat its installation step. After the fifth command, open a new PowerShell window and run
\texttt{pdftoppm -v}.
\begin{ps}[Repair of the search path, command 1 of 5]
cd "$env:LOCALAPPDATA\Microsoft\WinGet\Packages"
\end{ps}
\begin{ps}[Repair of the search path, command 2 of 5]
$p = Get-ChildItem -Recurse -Filter pdftoppm.exe
\end{ps}
\begin{ps}[Repair of the search path, command 3 of 5]
$p[0].DirectoryName
\end{ps}
\begin{ps}[Repair of the search path, command 4 of 5]
$old = [Environment]::GetEnvironmentVariable("Path", "User")
\end{ps}
\begin{pswide}[Repair of the search path, command 5 of 5]
[Environment]::SetEnvironmentVariable("Path", "$old;$($p[0].DirectoryName)", "User")
\end{pswide}

\trouble{\texttt{Python was not found; run without arguments to install from the Microsoft
Store}, or the Microsoft Store opens when a command beginning with \texttt{python} is run}
\cause The environment is not activated in this window, so Windows offers its own placeholder for
Python.
\remedy Check the prompt. If it does not begin with \texttt{(.venv)}, give the opening sequence
(the two commands at the head of this section).

\trouble{\texttt{py : The term 'py' is not recognized}, or \texttt{Python 3.12 not found}}
\cause Python 3.12 is not installed, or the window is older than the installation.
\remedy Open a new window. If the message persists, repeat the Python installation step of
Section~\ref{sec:tools}.

\trouble{\texttt{Activate.ps1 cannot be loaded because running scripts is disabled on this
system}}
\cause PowerShell's execution policy forbids scripts.
\remedy Run Step~\ref{st:policy} and answer \texttt{Y}. If the policy cannot be changed because the
laptop is managed by an organization, the environment can be used without activation: in every
command of this guide and its companion, write \code|.\.venv\Scripts\python.exe| in place of the
word \texttt{python} at the beginning, and \code|.\.venv\Scripts\python.exe -m pip| in place of
\texttt{pip}.

\trouble{\texttt{The module '.venv' could not be loaded}, with the activation line shown
twice in a row}
\cause The activation command was pasted twice before \key{Enter} was pressed.
\remedy Press \key{Esc} to clear the line, then paste the command once.

\trouble{\texttt{.\bs .venv\bs Scripts\bs Activate.ps1 : The term \ldots{} is not recognized}}
\cause The current folder is not the project folder, so there is no \file{.venv} in it.
\remedy Run \code|cd C:\ee105\camera-to-latex| first: the opening sequence has two commands, and the
order matters.

\trouble{\texttt{The '<' operator is reserved for future use}}
\cause A command copied from another source contained a placeholder in angle brackets, such as
\texttt{<your-username>}, and it was pasted without being replaced.
\remedy Use the commands of this guide, which contain no such placeholders.

\trouble{The window shows no new output and does not react}
\cause If the title bar contains the word \emph{Select}, a mouse click has paused the program.
Otherwise the command is probably still working.
\remedy Press \key{Esc} in the first case. In the second, wait; Steps~\ref{st:pipinstall}
and~\ref{st:train} state how long the slow commands may take.

\subsection{Python packages}

\trouble{\texttt{ERROR: Could not find a version that satisfies the requirement tensorflow}}
\cause The environment was created with a Python version for which TensorFlow is not published.
\remedy Run \code|python --version| with the environment activated. If it does not print
\texttt{Python 3.12}, close the window, delete the folder \file{.venv} in File Explorer, and
repeat Section~\ref{sec:python} from its second step, typing \texttt{py -3.12} exactly.

\trouble{\texttt{OSError: [Errno 2] No such file or directory} with the hint \texttt{Windows
Long Path support}}
\cause The project folder lies so deep in the folder tree that some file names of TensorFlow exceed
260 characters.
\remedy Install the project in \file{C:\bs ee105} as described in Section~\ref{sec:clone}, not in a
folder inside \file{Desktop} or \file{Documents}.

\trouble{\texttt{ReadTimeoutError}, or \texttt{Connection \ldots{} timed out}, during
\texttt{pip install}}
\cause The internet connection was interrupted.
\remedy Run the same command again. Packages already downloaded are not fetched a second time.

\trouble{\texttt{AttributeError: module 'cv2' has no attribute 'HOGDescriptor'}}
\cause Version~5 of OpenCV is installed.
\remedy Run Step~\ref{st:opencv}.

\trouble{\texttt{ModuleNotFoundError: No module named 'cv2'} (or \texttt{numpy},
\texttt{tensorflow})}
\cause The environment is not activated, or Step~\ref{st:pipinstall} did not finish.
\remedy Check for \texttt{(.venv)} in the prompt; if it is present, run Step~\ref{st:pipinstall}
again.

\trouble{\texttt{No module named pipeline} (or \texttt{tools}, \texttt{model})}
\cause The command was run outside the project folder.
\remedy Run \code|cd C:\ee105\camera-to-latex| and repeat the command.

\trouble{\texttt{ImportError: DLL load failed while importing \_pywrap\_tensorflow\_internal}}
\cause A Microsoft runtime library that TensorFlow requires is missing from Windows.
\remedy Run \code|winget install --id Microsoft.VCRedist.2015+.x64 -e --source winget|, open a new
window, give the opening sequence (the two commands at the head of this section), and repeat
Step~\ref{st:verifypy}.

\trouble{Lines that mention \texttt{oneDNN custom operations}, \texttt{absl::InitializeLog}
or \texttt{GPU support is not available on native Windows}}
\cause TensorFlow announces how it will compute. On Windows it uses the processor, not a graphics
card, which is sufficient for this project.
\remedy None is needed; these are not errors.

\subsection{\LaTeX{} and Poppler}

\trouble{\texttt{'pdflatex' is not on PATH}, or \texttt{'pdflatex' not found}}
\cause MiKTeX is not installed, or the window is older than its installation.
\remedy Open a new window. If the message persists, repeat the MiKTeX steps of
Section~\ref{sec:tools}.

\trouble{\texttt{pdflatex took longer than 20 s}}
\cause MiKTeX was downloading a component, or was waiting behind a dialog for permission to do so.
\remedy Look for a MiKTeX dialog behind the other windows and click \ui{Install}. Run
Step~\ref{st:autoinstall}, then repeat the command; the second attempt is fast because the
component is now present.

\trouble{\texttt{pdflatex failed:} followed by \texttt{LaTeX Error: File \ldots{} not found}}
\cause MiKTeX could not fetch a component, usually because it has never synchronized with its
servers or because there is no internet connection.
\remedy Open MiKTeX Console, click \ui{Check for updates}, install what is offered, and repeat the
command.

\trouble{\texttt{RuntimeError: pdftoppm (poppler) is needed for -{}-render}}
\cause Poppler is not installed or not found.
\remedy Run \texttt{pdftoppm -v} in a new window and follow the first entries of this section.

\subsection{Commands of the project}

\trouble{\texttt{symbols\_v1.json not found: train a model first}}
\cause No classifier has been trained on this laptop yet.
\remedy Run Steps~\ref{st:synth} and~\ref{st:train}.

\trouble{\texttt{missing \ldots{}train.npz: run python -m model.make\_synthetic first}}
\cause Training was started before the training images were generated.
\remedy Run Step~\ref{st:synth}, then Step~\ref{st:train}.

\trouble{\texttt{no photos found in test\_set\bs images for split test}}
\cause The evaluation was run without its two options, so it looked for photographs taken with the
board, which do not exist yet.
\remedy Use the full command of Step~\ref{st:eval}, which ends in
\code|--images test_set\synthetic --split all|.

\trouble{\texttt{no results file: run python -m tools.eval\_pipeline first}}
\cause The review document is built from an evaluation, and none has been run.
\remedy Run Step~\ref{st:eval}, then Step~\ref{st:review}.

\trouble{\texttt{camera problem:} followed by a remark about ports}
\cause The application was started without the option \texttt{-{}-image} and looked for the board.
\remedy Use the full command of Step~\ref{st:onephoto}. Work with the board is described in the
companion document.

\subsection{Git}

\trouble{\texttt{remote: Repository not found.}}
\cause Your GitHub account has no access to the repository, or the wrong account is signed in.
\remedy See the amber box of Step~\ref{st:gitclone}.

\trouble{\texttt{fatal: destination path 'camera-to-latex' already exists and is not an
empty directory}}
\cause The repository was cloned earlier.
\remedy Nothing; continue with the step after the clone.

\trouble{\texttt{Author identity unknown} \ldots{} \texttt{Please tell me who you are}}
\cause Git does not know your name and email address.
\remedy Carry out the subsection \emph{Introduce yourself to Git} of Section~\ref{sec:clone}.

\trouble{\texttt{fatal: not a git repository}}
\cause The current folder is not the project folder.
\remedy Run \code|cd C:\ee105\camera-to-latex|.

\trouble{\texttt{[rejected] main -> main}, \texttt{CONFLICT}, or \texttt{Your local changes
\ldots{} would be overwritten}}
\cause The team's copy and yours have diverged.
\remedy Section~\ref{sec:routine} treats each of the three cases.

\subsection{The laptop}

\trouble{An installation inside \file{Desktop} or \file{Documents} is extremely slow}
\cause OneDrive synchronizes those folders and uploads every file of the environment.
\remedy Keep the project in \file{C:\bs ee105}. If it was cloned elsewhere, delete that copy and
repeat Sections~\ref{sec:clone} and~\ref{sec:python}.

\trouble{Training stopped because the laptop went to sleep}
\cause The power settings suspended the computation.
\remedy Switch sleep off while the charger is connected, as described before Step~\ref{st:train}, and
start the training again.
